\chapter{Introduction}
\label{ch:introduction}


\section{Background and motivation}
\label{sec:background}


An accelerated global development of reverse osmosis (RO) membranes has created a rapidly growing volume of end-of-life (EoL) membrane waste that is difficult to manage through existing disposal pathways. As membrane-based desalination and water treatment gains increasing investment, the quantity of membranes requiring replacement has also grown proportionally. However, most membranes continue to be disposed of to landfill or incineration, and this presents a missed opportunity for material recovery and reuse.


\section{Problem statement}
\label{sec:problem}
Names shortcoming in current practic

\begin{itemize}
  \item EoL membranes are landfilled or incinerated.
  \item Downcycling (RO$\rightarrow$NF/UF) recovers a lower-grade product and permanently discards RO selectivity.
  \item Regenerating the polyamide layer would restore that selectivity, but has been demonstrated only in a handful of studies, generally on pretreated substrates
\end{itemize}


\section{Research aim}
\label{sec:aim}
This research aims to evaluate whether EoL RO membranes can be uplifted by depositing a brand new polyamide layer via ex-situ interfacial polymerisation, and how the fouling present and cleaning pre-treatement affects the membrane's performance in permeability and salt rejection

\section{Objectives}
\label{sec:objectives}

\begin{enumerate}
  \item Characterise the EoL RO membranes and their fouling profiles.
  \item Prepare substrates --- uncleaned and extensively cleaned.
  \item Regenerate the polyamide layer by ex-situ interfacial polymerisation.
  \item Measure pure water permeability and NaCl rejection across two membranes.
  \item Autopsy selected coupons to relate surface condition to performance.
\end{enumerate}




\section{Scope and limitations}
\label{sec:scope}
Defines the boundaries of the study so that the reader does not expect results the experimental design cannot deliver.

\begin{itemize}
  \item In scope --- fouling characterisation of membranes A and C; two cleaning arms (none, extensive); in-situ IP; PWP and NaCl rejection; autopsy of selected coupons.
  \item Out of scope --- optimisation of IP or cleaning chemistry; life cycle assessment; techno-economic analysis; module-scale or long-duration trials.
  \item Constraints --- limited replicates, single compaction protocol, single feed chemistry.
\end{itemize}



\section{Thesis outline}
\label{sec:outline}
Maps the remaining chapters

% One paragraph: Ch.2 reviews the literature and establishes the gap;
% Ch.3 describes materials and methods; Ch.4 presents and discusses
% results; Ch.5 concludes and recommends further work.
